%% problem_id: S014-lurie-ultracategories-prop-3.4.6
%% source_id: S014 (Jacob Lurie, "Ultracategories")
%% source_url: https://www.math.ias.edu/~lurie/papers/Conceptual.pdf
%% source_file: lurie-conceptual.pdf
%% statement_locator: printed p. 36, Proposition 3.4.6
%% proof_locator: printed pp. 38--41, section 3.5 "The Proof of Proposition 3.4.6"
%%   (Lemma 3.5.1 p. 38; Notation 3.5.2 and the construction of u, v pp. 39--40;
%%    Lemmas 3.5.3 and 3.5.4 p. 40; Lemma 3.5.5 and the final assembly p. 41)
%% representation: PDF; transcribed by RENDERING the pages (pdftoppm -r 150 -png) and READING THE
%%   IMAGES, then checking every display against them.  (The paper's text layer mangles symbols, so it
%%   was used only for locating pages, never as the text of record.)
%% difficulty_level: H3
%% why H3 (and why NOT Theorem 3.4.4): the paper says of Theorem 3.4.4 that "the main ingredient in our
%%   proof of Theorem 3.4.4 is the following result, whose proof we defer to section 3.5" -- i.e. 3.4.4
%%   delegates its decisive step.  Section 3.5 is where that step IS carried out: the surjectivity of
%%   theta: F_{G,y} -> G(y), proved by a four-page self-contained construction (a two-stage
%%   Stone-Cech/ultraproduct setup, the categorical Fubini transformation, and an equalizer
%%   identification), with Lemmas 3.5.1/3.5.3/3.5.4/3.5.5 all proved inside the section.  So the
%%   collection rule "收集承担者而不是招牌" applied and this is the carrier.
%% status: complete (statement + author's proof verbatim + reason + locators)
%% ---------------------------------------------------------------------------
\documentclass[11pt]{article}
\usepackage{amsmath,amssymb,amsthm}
\newtheorem{theorem}{Theorem}
\newtheorem{lemma}{Lemma}
\newtheorem{proposition}{Proposition}
\newtheorem{notation}{Notation}
\newcommand{\Ult}{\mathrm{LUlt}}
\begin{document}

\section*{Proposition 3.4.6}

\textbf{Source.} J.~Lurie, \emph{Ultracategories}; printed page 36 (statement) and pages 38--41
(section 3.5, the proof).

\begin{proposition}
Let $X$ be a compact Hausdorff space and let $G: X \to \mathrm{Set}$ be a left ultrafunctor. Then,
for each point $y \in X$, the preceding construction induces a bijection
\[
\mathscr F_{G,y} \longrightarrow G(y).
\]
\end{proposition}

\subsection*{Proof (authors' text, section 3.5)}

%% --- page 38 ---
Throughout this section, we fix a compact Hausdorff space $X$ and a left ultrafunctor
$G: X \to \mathrm{Set}$. Our goal is to compute the stalks of the sheaf $\mathscr F_G$ of
Construction 3.4.2. We begin by establishing a weak version of Proposition 3.4.6.

\begin{lemma}[3.5.1]
For each point $y \in X$, the map
$\mathscr F_{G,y} = \varinjlim_{y \in U} \mathscr F_G(U) \to G(y)$ of Proposition 3.4.6 is injective.
\end{lemma}

\begin{proof}
Let $U$ be an open neighborhood of the point $y$ and suppose we are given a pair of elements
$\phi_\bullet = \{\phi_x\}_{x \in U}$ and $\psi_\bullet = \{\psi_x\}_{x \in U}$ of $\mathscr F_G(U)$
satisfying $\phi_y = \psi_y$ (as elements of the set $G(y)$). We wish to show that there exists an open
subset $V \subseteq U$ containing the point $y$ such that $\phi_x = \psi_x$ for each $x \in V$.
Assume otherwise. Set $U_0 = \{x \in U : \phi_x \neq \psi_x\} \subseteq U$, and let $\mathcal U$ be the
collection of all subsets of $U$ having the form $U_0 \cap V$, where $V$ is an open neighborhood of $y$.
Then $\mathcal U$ is closed under finite intersections and does not contain the empty set. Applying
Proposition 1.1.10 we deduce that there exists an ultrafilter $\mu$ on the set $U$ such that
$\mu(U_0 \cap V) = 1$, whenever $V$ is an open neighborhood of $y$. We then have $\int_U x\,d\mu = y$.
Invoking our assumption that $\phi_y = \psi_y$ and the definition of the sheaf $\mathscr F_G$, we deduce
that $\phi_\bullet$ and $\psi_\bullet$ have the same image in the ultraproduct $\int_U G(x)d\mu$. It
follows that $\mu(U_0) = 0$, contradicting the definition of $U_0$.
\end{proof}

%% --- page 39 ---
To complete the proof of Proposition 3.4.6 we must show that each element $\varphi_y \in G(y)$ can be
extended to an element $\phi_\bullet \in \mathscr F_G(U)$ for some open neighborhood $U$ of the point
$y$. The proof will require some auxiliary constructions.

\begin{notation}[3.5.2]
Choose a set $T$ equipped with a bijection $f_0: T \to X$. We regard $f_0$ as a continuous bijection of
topological spaces, where $T$ is endowed with the discrete topology. We will use the letter $\nu$ to
denote a typical ultrafilter on $T$: that is, a point in the Stone-\v{C}ech compactification $\beta T$.
Let $i: T \hookrightarrow \beta T$ be the canonical embedding, which assigns to each point $t \in T$ the
corresponding principal ultrafilter $i(t) = \delta_t$. By virtue of Proposition 3.2.7 the map
$f_0: T \to X$ extends uniquely to a continuous map $f: \beta T \to X$, given concretely by the formula
$f(\nu) = \int_T f_0(t)d\nu$.

Let $\mathscr G_0$ denote the sheaf of sets on the (discrete) topological space $T$ whose stalks are
given by $\mathscr G_{0t} = G(f_0(t))$, and let $\mathscr G$ denote the direct image $i_* \mathscr G_0$.
Then $\mathscr G$ is a sheaf of sets on $\beta T$ whose value on closed and open sets is given by the
formula
\[
\mathscr G(\beta T_0) = \prod_{t \in T_0} G(f_0(t)),
\]
and whose stalk at a point $\nu \in \beta T$ is given by the formula
$\mathscr G_\nu \simeq \int_T G(f_0(t))d\nu$.

Choose a set $S$ and a bijection $g_0: S \to \beta T$, which we identify with a collection of
ultrafilters $\{\nu_s\}_{s \in S}$ on the set $T$. We regard $g_0$ as a continuous bijection of
topological spaces, where $S$ is equipped with the discrete topology and $\beta T$ with the topology of
Construction 3.2.2. We will use the letter $\mu$ to denote a typical ultrafilter on $S$: that is, a
point in the Stone-\v{C}ech compactification $\beta S$. Let $j: S \hookrightarrow \beta S$ be the
canonical embedding, which assigns to each point $s \in S$ the corresponding principal ultrafilter
$j(s) = \delta_s$. Let $\mathscr H_0 = g_0^* \mathscr G$ denote the sheaf of sets on $S$ whose stalks are
given by $\mathscr H_{0s} = \int_T G(f_0(t))d\nu_s$. Let $\mathscr H$ denote the sheaf of sets on
$\beta S$ given by the direct image $j_* \mathscr H_0$. Then $\mathscr H$ is given on closed and open
sets by the formula
\[
\mathscr H(\beta S_0) = \prod_{s \in S_0} \int_T G(f_0(t))d\nu_s,
\]
and its stalk at a point $\mu \in \beta S$ is given by
$\mathscr H_\mu \simeq \int_S \bigl(\int_T G(f_0(t))d\nu_s\bigr)d\mu$.

By virtue of Proposition 3.2.7 the bijection $g_0: S \to \beta T$ admits a unique continuous extension
$g: \beta S \to \beta T$, given concretely by the formula $g(\mu) = \int_S \nu_s d\mu$ (see Example
3.2.8). We have a canonical isomorphism $j^* g^* \mathscr G \simeq g_0^* \mathscr G = \mathscr H_0$ which
induces a map of sheaves $u: g^* \mathscr G \to j_* \mathscr H_0 = \mathscr H$. Concretely, the map $u$
is given on stalks by the map
\[
(g^* \mathscr G)_\mu = \mathscr G_{g(\mu)} = \int_T G(f_0(t))d\Bigl(\int_S \nu_s d\mu\Bigr)
\xrightarrow{\ \Delta_{\lambda,\mu,*} \ }
\int_S \Bigl(\int_T G(f_0(t))d\nu_s\Bigr)d\mu \simeq \mathscr H_\mu,
\]
where $\Delta_{\lambda,\mu,*}$ is the categorical Fubini transformation in the category of sets.

Let $h_0: S \to T$ be the map of sets given by the composition
\[
S \xrightarrow{\ g_0\ } \beta T \xrightarrow{\ f\ } X \xrightarrow{\ f_0^{-1}\ } T \xrightarrow{\ i\ } \beta T .
\]
For each $s \in S$, the right ultrastructure on the functor $G$ determines a map
\[
\mathscr G_{(i \circ h_0)(s)} \simeq \mathscr G_{0h_0(s)} \simeq G((f \circ g_0)(s))
= G\Bigl(\int_T f_0(t)d\nu_s\Bigr) \xrightarrow{\ \sigma_{\nu_s}\ }
\int_T G(f_0(t))d\nu_s \simeq \mathscr H_{0s}.
\]
Let $h: \beta S \to \beta T$ denote the unique continuous extension of $h_0$, given concretely by
$h(\mu) = h_{0*}(\mu)$. The preceding discussion then gives a map of sheaves
$j^* h^* \mathscr G \to \mathscr H_0$ on the discrete space $S$, which we identify with a map
$v: h^* \mathscr G \to j_* \mathscr H_0 = \mathscr H$ on the topological space $\beta S$.

%% --- page 40 ---
Concretely, the map $v$ is given on stalks by the composition
\[
\begin{aligned}
(h^* \mathscr G)_\mu &\simeq \int_T G(f_0(t))d(h_{0*}\mu) \\
&\xrightarrow{\ \Delta_{\mu,h_0}\ } \int_S G((f_0 \circ h_0)(s))d\mu \\
&= \int_S G\Bigl(\int_T f_0(t)d\nu_s\Bigr)d\mu \\
&\xrightarrow{\ \int_S \sigma_{\nu_s} d\mu\ } \int_S \Bigl(\int_T G(f_0(t))d\nu_s\Bigr)d\mu .
\end{aligned}
\]
Note that we have a commutative diagram of topological spaces
\[
\beta S \underset{h}{\overset{g}{\rightrightarrows}} \beta T \xrightarrow{\ f\ } X ;
\]
that is, we have $f \circ g = f' = f \circ h$ for some continuous map $f': \beta S \to X$, given
concretely by the formula $f'(\mu) = \int_S \bigl(\int_T f_0(t)d\nu_s\bigr)d\mu$. The maps $u$ and $v$
therefore determine maps $u', v': f_*(\mathscr G) \to f'_*(\mathscr H)$ in the category
$\mathrm{Shv}(X)$.

\begin{lemma}[3.5.3]
The sheaf $\mathscr F_G$ of Construction 3.4.2 can be identified with the equalizer of the maps
$u', v': f_*(\mathscr G) \rightrightarrows f'_*(\mathscr H)$.
\end{lemma}

\begin{proof}
By construction, the direct image sheaf $f_* \mathscr G \simeq f_{0*} \mathscr G_0$ is given by
\[
(f_* \mathscr G)(U) \simeq \prod_{x \in U} G(x).
\]
Let $\mathscr E$ denote the equalizer of $u'$ and $v'$, regarded as a subsheaf of $f_* \mathscr G$.
Unwinding the definitions, we see that for each subset $U \subseteq X$, we can identify
$\mathscr E(U)$ with the subset of $\prod_{x \in U} G(x)$ consisting of those tuples
$\{\phi_x\}_{x \in U}$ which satisfy the following condition:
\begin{itemize}
\item[(★′)] For every ultrafilter $\nu$ on $T$ such that $\int_T f_0(t)d\nu$ belongs to $U$, the elements
$\{\phi_{f_0(t)}\}_{t \in f_0^{-1}(U)}$ and $\phi_{\int_T f_0(t)d\nu}$ have the same image under the maps
\[
\prod_{t \in f_0^{-1}(U)} G(f_0(t)) \xrightarrow{\ q_\nu^{f_0^{-1}(U)}\ } \int_T G(f_0(t))d\nu
\xleftarrow{\ \sigma_\nu\ } G\Bigl(\int_T f_0(t)d\nu\Bigr),
\]
where $\sigma_\nu$ is given by the left ultrastructure on $G$.
\end{itemize}
Note that condition (★′) follows immediately from condition (∗) of Construction 3.4.2 (applied to the
subset $T_0 = f^{-1}(U) \subseteq T$, and the ultrafilter on $T_0$ given by the restriction of $\nu$).
That is, we can regard the sheaf $\mathscr F_G$ of Construction 3.4.2 as a subsheaf of $\mathscr E$. To
complete the proof, we must show that (★′) implies (∗). To this end, suppose we are given an arbitrary
map of sets $e: R \to U$ and an ultrafilter $\lambda$ on $R$ satisfying $\int_R e(r)d\lambda \in U$; we
wish to show that $\{\phi_{e(r)}\}_{r \in R}$ and $\phi_{\int_R e(r)d\lambda}$ have the same image under
the maps
\[
\prod_{r \in R} G(e(r)) \xrightarrow{\ q_\lambda\ } \int_R G(e(r))d\lambda
\xleftarrow{\ \sigma_\lambda\ } G\Bigl(\int_R e(r)d\lambda\Bigr).
\]
Since $f_0$ is bijective, the map $e$ factors uniquely as a composition
$R \xrightarrow{\ e'\ } T_0 \xrightarrow{\ f_0\ } X$. The desired assertion now follows by applying (∗)
to the ultrafilter $e'_* \lambda$ on $T_0 = f_0^{-1}(U)$.
\end{proof}

Let us now fix a point $y \in X$ and an element $\phi_y$ of the set $G(y)$. For each ultrafilter $\nu$ on
$T$ satisfying $f(\nu) = \int_T f_0(t)d\nu = y$, let $\psi_\nu$ denote the image of $\phi_y$ under the map
\[
G(y) = G\Bigl(\int_T f_0(t)d\nu\Bigr) \xrightarrow{\ \sigma_\mu\ } \int_T G(f_0(t))d\nu = \mathscr G_\nu .
\]

\begin{lemma}[3.5.4]
Let $\mu$ be a point of $\beta S$ satisfying $f'(\mu) = y$. Let
$u_\mu: \mathscr G_{g(\mu)} \to \mathscr H_\mu$ and $v_\mu: \mathscr G_{h(\mu)} \to \mathscr H_\mu$ be
the maps induced by $u$ and $v$, respectively. Then
$u_\mu(\psi_{g(\mu)}) = v_\mu(\psi_{h(\mu)})$.
\end{lemma}

%% --- page 41 ---
\begin{proof}
This follows from the definition of $u$ and $v$, together with the commutativity of the diagram
\[
\begin{array}{ccc}
G(y) & \xrightarrow{\ \sigma_{h_{0*}\mu}\ } & \int_T G(f_0(t))d(h_{0*}\mu) \\
\Big\Vert & & \Big\downarrow{\scriptstyle \Delta_{\mu,h_0}} \\
G\bigl(\int_T f_0(t)d(\int_S \nu_s d\mu)\bigr) & & \int_S G((f_0 \circ h_0)(s))d\mu \\
\Big\downarrow{\scriptstyle \sigma_{\int_S \nu_s d\mu}} & & \Big\downarrow{\scriptstyle \int_S \sigma_{\nu_s} d\mu} \\
\int_T G(f_0(t))d\bigl(\int_S \nu_s d\mu\bigr) & \xrightarrow{\ \Delta_{\mu,\nu_*}\ } &
\int_S \bigl(\int_T G(f_0(t))d\nu_s\bigr)d\mu .
\end{array}
\]
\end{proof}

\begin{lemma}[3.5.5]
There exists a global section $\psi$ of the sheaf $\mathscr G|_{f^{-1}\{y\}}$ such that, for each point
$\nu \in f^{-1}\{y\}$, the germ of $\psi$ at $\nu$ is equal to $\psi_\nu$.
\end{lemma}

\begin{proof}
Let $\overline y$ denote the unique element of $T$ satisfying $f_0(\overline y) = y$, and let
$K \subseteq \beta S$ denote the inverse image of $\delta_{\overline y}$ under the map
$h: \beta S \to \beta T$. Then $g: \beta S \to \beta T$ restricts to a map of topological spaces
$g_K: K \to f^{-1}\{y\}$. Then $u$ restricts to a map of sheaves
$u_K: g_K^* \mathscr G|_{f^{-1}\{y\}} \to \mathscr H|_K$ on $K$, which we can identify with a map
$\iota: \mathscr G|_{f^{-1}\{y\}} \to (u_{K*} \mathscr H|_K)$ on the fiber $f^{-1}\{y\}$. We first claim
that $\iota$ is injective. To prove this, it suffices to observe that for each point
$\nu \in f^{-1}\{y\}$, there exists a point $\mu \in K$ for which $u_K(\mu) = \nu$ and $u$ induces a
monomorphism of stalks $u_\mu: \mathscr G_\mu \to \mathscr H_\nu$. In fact, writing
$\nu = \nu_s$ for some $s \in S$, we can take $\mu = \delta_s$ to be the principal ultrafilter at the
point $s$, in which case the map $u_\mu$ is bijective.

For each point $\nu \in f^{-1}\{y\}$, let $\psi'_\nu$ denote the image of $\psi_\nu$ in the stalk
$(u_{K*} \mathscr H|_K)_\nu$ (under the map $\iota$). Since $\iota$ is a monomorphism, it will suffice to
show that the germs $\{\psi'_\nu\}_{\nu \in f^{-1}\{y\}}$ determine a global section of the sheaf
$u_{K*} \mathscr H|_K$. Applying Proposition 0.0.10 to the map $u_K: K \to f^{-1}\{y\}$, this can be
reformulated as follows:
\begin{itemize}
\item[(★)] There exists a global section $\xi$ of the sheaf $\mathscr H|_K$ having the property that, for
every point $\mu \in K$, the germ $\xi_\mu$ of $\xi$ at the point $\mu$ coincides with the image of
$\psi_{g(\mu)}$ under the map $\mathscr G_{g(\mu)} \to \mathscr H_\mu$ determined by $u_K$.
\end{itemize}
Let $\underline{G(y)}_K$ denote the constant sheaf on $K$ with the value $G(y)$, so that $v$ restricts
to a map of sheaves
\[
v_K: \underline{G(y)}_K \to \mathscr H|_K .
\]
The map $v_K$ carries the element $\phi_y \in G(y)$ to a global section $\xi$ the sheaf
$\mathscr H|_K$. It follows from Lemma 3.5.4 that $\xi$ satisfies the requirement of (★).
\end{proof}

\subsection*{Proof of Proposition 3.4.6 (authors' text)}

Let $\mathscr G(f^{-1}\{y\})$ denote the set of global sections of the sheaf
$\mathscr G|_{f^{-1}\{y\}}$, and define $\mathscr H(f^{-1}\{y\})$ similarly. Since $f$ and $f'$ are
proper maps, we can identify $\mathscr G(f^{-1}\{y\})$ and $\mathscr H(f^{-1}\{y\})$ with the stalks
$(f_* \mathscr G)_y$ and $(f'_* \mathscr H)_y$, respectively. Using Lemma 3.5.3 we obtain an equalizer
diagram of sets
\[
\mathscr F_{G,y} \to \mathscr G(f^{-1}\{y\}) \rightrightarrows \mathscr H(f^{-1}\{y\}).
\]
For each element $\phi_y \in G(y)$, the section $\psi \in \mathscr G(f^{-1}\{y\})$ of Lemma 3.5.5
belongs to this equalizer (by Lemma 3.5.4), and can therefore be identified with an element of
$\mathscr F_{G,y}$. It follows immediately from the construction that the canonical map
$\theta: \mathscr F_{G,y} \to G(y)$ carries this element to $\phi_y$, which shows that $\theta$ is
surjective; injectivity was established in Lemma 3.5.1. \hfill$\square$

\end{document}
